\chapter{统计状态的微观建模：信息几何、估计界与学习坐标}
\label{chp:microscopic_origins}

我在本章讨论的“微观”，不是把智能还原成某种物理粒子，而是把一个可观察的智能状态继续分解为样本、概率、参数、估计量和更新记录。前章已经说明，表示距离只有在对象、度量与任务商空间确定之后才有意义；本章进一步追问这些局部度量可以从什么统计问题中得到，又在哪些条件下失效。信息几何在这里是一套分析统计模型可辨识性与更新灵敏度的工具，而不是关于世界本体的证明。我们将分别处理模型族、Fisher 信息、估计下界、分布表示和学习坐标，并把“定义”“正则模型中的数学结果”“有限数据下的经验判断”与“AgentOS 的工程选择”逐层分开。这样，HSF-HD 所说的一般状态动力学可以吸收统计学习的可靠结构，同时不把某一概率参数化、某一种优化器或某一种硬件机制规定为一切智能的共同微观实体。

\section{从任务观测到统计模型：世界不等于分布族}
\label{sec:section_8204}

我先固定建模对象。设 $\mathcal X$ 是一次任务观测的样本空间，$\mathcal A$ 是其事件族，$\nu$ 是选定的参考测度；当观测离散时，$\nu$ 可以是计数测度，当观测连续时，它可以是与坐标约定相容的测度。参数空间记为 $\Theta\subseteq\mathbb R^d$，其中 $d$ 是有限的参数维数。统计模型不是“整个世界”，而是相对于观测协议形成的分布族
\begin{equation}
    \mathcal P=\{P_\theta:\theta\in\Theta\},\qquad
    p_\theta(x)=\frac{dP_\theta}{d\nu}(x),\qquad
    \int_{\mathcal X}p_\theta(x)\,d\nu(x)=1.
\end{equation}
这里 $P_\theta$ 的单位是无量纲概率；若写密度 $p_\theta$，它相对于 $\nu$ 的量纲必须使积分无量纲。真实数据机制记为 $Q$，它可以在模型族之外。于是“选择参数”只是用 $P_\theta$ 近似任务相关的 $Q$，并不声称 $Q$ 包含环境的全部状态，更不声称环境本身由参数坐标生成。摄像头采样、日志事件和语言片段即使指向同一外部对象，也会因为采样率、遮挡规则、分词方案和记录权限不同而产生不同的 $\mathcal X$、$Q$ 与可辨识结构。

这一区分使我能够准确解释原章所谓“世界图即分布族”。HSF-HD 的世界结构快照是带身份、类型、来源和版本的关系记录，可写成 $G_t=(V_t,E_t,\tau_t,\sigma_t)$；概率模型则描述在给定快照、观测接口与任务条件后可能出现什么记录。若以 $Y$ 表示任务读出，以 $Z$ 表示局部结构身份，以 $X$ 表示当前观测，一个可能的工程分解是
\begin{equation}
    P_\theta(Y,Z,X\mid G_t,c_t)
    =P_\theta(Y\mid Z,X,G_t,c_t)
     P_\theta(Z\mid X,G_t,c_t)
     P_\theta(X\mid G_t,c_t),
\end{equation}
其中 $c_t$ 汇集目标、边界和观测协议。这个分解是一种模型假设，不是概率乘法自动给出的因果顺序。它允许局部结构网络 $G_t$ 与全局分布表示同时出现：前者保存哪些节点和关系可被寻址，后者在大量候选上表达不确定性和相似性；二者通过任务相关身份 $Z$、绑定条件及读出 $Y$ 联合，而不能由“脑区”或“流形”的相似图像推出机制等同。

参数也不应被直接解释为语义对象。令 $R_\theta:\mathcal X\to\mathbb R^m$ 是分布式编码，$D_\theta$ 是从编码与结构记录到任务输出的读出。两个参数 $\theta$ 和 $\theta'$ 若对所有允许输入都给出相同的条件分布与读出，就在当前任务下不可区分。定义等价关系
\begin{equation}
    \theta\sim\theta'
    \Longleftrightarrow
    P_\theta(Y\in B\mid X=x,G,c)
    =P_{\theta'}(Y\in B\mid X=x,G,c)
\end{equation}
对所有合法的 $x,G,c$ 和可测集合 $B$ 成立，则真正可辨识的状态位于商空间 $\Theta/{\sim}$，而不是原始坐标表。神经网络中的通道置换、隐层尺度补偿或矩阵分解重参数化，常使多个参数点实现同一函数。若忽略这一点，某些方向的 Fisher 信息退化并不表示“概念没有意义”，只说明现有观测与读出不能区分那些参数变化。反过来，参数数值相近也不保证结构身份相同，因为一次离散的节点合并、来源撤销或权限变更可能在小参数扰动之外改变任务语义。

接下来需要说明坐标变化的条件。若 $\phi=f(\theta)$ 是同一可辨识模型区域内的光滑双射，且雅可比 $J=\partial\theta/\partial\phi$ 可逆，那么 $p_\theta=p_{\theta(\phi)}$ 只是重参数化；由此得到的局部统计量必须按坐标变换规则共同改变。若 $f$ 不可逆、模型支持集发生变化、维数改变，或者结构图执行了离散编辑，就不是同一坐标系之间的改名，不能沿用固定维数流形的微分结论。实际系统还必须记录模型版本 $v$、数据时间窗 $I$、抽样策略 $s$ 和任务上下文 $c$，所以一个可审计统计状态至少应写成 $(\theta,v,I,s,c)$。只保留向量而丢弃这些索引，会把数据漂移、版本替换和真实状态变化混成同一种“几何移动”。

最后，我把这一节的主张限制在它能够承担的范围内。概率族提供了表示不确定性、比较局部分布与分析估计误差的语言，但它不保证模型正确，也不自动给出因果机制。有限样本拟合优良只能支持规定数据与指标上的经验结论，不能推出内部模型与世界整体同构。AgentOS 可以把当前模型参数、世界结构快照、外部记忆 $M_e$、证据来源和任务边界封装为版本化状态，再由 SPC、TCE 或其他调度组件选择更新；这是把上述条件落实为可审计接口的一种工程方案，而不是一般智能必须采用的唯一结构。后续 Fisher 度量和估计下界，都只在这一已声明的统计实验内推导。
